\documentclass[10pt]{article}
\usepackage[paperwidth=11in,paperheight=17in,margin=0.6in]{geometry}
\usepackage[T1]{fontenc}
\usepackage{mathpazo}
\usepackage{amsmath,amssymb,bm,mathtools}
\usepackage{booktabs,array,tabularx,multicol}
\usepackage{xcolor,colortbl}
\usepackage[skip=4pt plus 1pt]{parskip}
\usepackage{ragged2e}
\usepackage{enumitem}
\usepackage{hyperref}

\definecolor{lithioxbggreen}{HTML}{8ABE53}
\definecolor{lithioxfontblue}{HTML}{001F3F}
\definecolor{lithioxfontgreen}{HTML}{32B641}
\definecolor{formteal}{HTML}{0891B2}
\definecolor{reportgrey}{HTML}{374151}
\definecolor{greytext}{HTML}{6B7280}
\hypersetup{colorlinks=true,linkcolor=lithioxfontblue,urlcolor=formteal}
\pagestyle{empty}
\AtBeginDocument{\RaggedRight}
\makeatletter
\renewcommand\section{\@startsection{section}{1}{0pt}%
  {10pt plus 2pt minus 1pt}{4pt}%
  {\large\sffamily\bfseries\color{lithioxfontgreen}}}
\renewcommand\subsection{\@startsection{subsection}{2}{0pt}%
  {7pt plus 1pt minus 1pt}{3pt}%
  {\normalsize\sffamily\bfseries\color{lithioxfontblue}}}
\makeatother
\newcommand{\dd}{\,\mathrm d}
\newcommand{\tr}{\operatorname{tr}}
\newcommand{\pinv}{^{\#}}

\begin{document}

\noindent\colorbox{lithioxbggreen!15}{%
\hspace{10pt}\begin{minipage}{\dimexpr\textwidth-2\fboxsep-20pt}
\vspace{10pt}
{\LARGE\sffamily\bfseries\color{lithioxfontblue} Full Molecular Electrolyte Conductivity}
\vspace{4pt}\RaggedRight

This document specifies one computational system that maps an electrolyte recipe,
temperature, pressure, and molecular species data to one bulk scalar conductivity.
The model resolves the complete periodic molecular liquid. Ion pairs, aggregates,
solvent shells, cages, and concentration fluctuations are sampled configurations.
They are not configured populations or conductivity corrections.
\vspace{10pt}
\end{minipage}\hspace{10pt}}

\vspace{0.3cm}
\begin{multicols}{2}
\RaggedRight\raggedcolumns

\section{Complete input and output}

The physical input is
\[
\mathcal I=(\mathcal R,T,p,\mathcal S,\Theta).
\]
The recipe \(\mathcal R\) gives solvent volume ratios, salt molarities, and additive
mass fractions. The species data \(\mathcal S\) give molecular topology, isotope
masses, formal molecular charge, interaction sites, and parameters for the common
interaction equations. The shared parameter set \(\Theta\) is fixed over the supported
chemical domain.

The numerical arguments are finite-cell scale \(m\), equilibrium integration
resolution \(r\), memory resolution \(\ell\), and basis level \(n\). The complete
calculation is
\[
\boxed{\begin{aligned}
(\mathcal R,T,p,\mathcal S,\Theta)
&\longmapsto \{N_s^{(m)}\},V_m^\star,q\\
&\longmapsto U_\Theta,\rho_\Theta\\
&\longmapsto K_{\Theta,0}(q),D_\Theta(q),L_\Theta,P\\
&\longmapsto A_n,h_n,C^{\rm dir}\\
&\longmapsto \sigma_{m,r,\ell,n}
\longmapsto \sigma(\mathcal R,T,p).
\end{aligned}}
\]

\section{Neutral periodic molecular system}

Let \(x_s(\mathcal R)\) be the molecular amount of component \(s\) after converting
the recipe units and expanding salts into molecular ions. At cell scale \(m\), choose
nonnegative integers
\[
\{N_s^{(m)}\}
=\arg\min_{\{n_s\}\subset\mathbb N_0}
\sum_s w_s\left(\frac{n_s}{m}-x_s(\mathcal R)\right)^2
\]
subject to
\[
\sum_s n_sQ_s=0,\qquad \mathbf Bn=\mathbf b_m.
\]
The first constraint imposes exact molecular electroneutrality. The matrix
\(\mathbf B\) contains exact stoichiometric constraints. The sequence obeys
\[
N_s^{(m)}/m\longrightarrow x_s(\mathcal R).
\]

For periodic cell matrix \(H\), volume \(V=\det H\), and constrained molecular
configuration space \(\Omega_V\), define
\[
Z_V(T,\{N_s\})=\int_{\Omega_V}
e^{-\beta U_\Theta(q;H)}\dd q,\qquad \beta=(k_BT)^{-1}.
\]
At pressure \(p\),
\[
\Delta(T,p,\{N_s\})=\int_0^\infty
e^{-\beta pV}Z_V(T,\{N_s\})\dd V,
\]
\[
V_m^\star=\Delta^{-1}\int_0^\infty
V e^{-\beta pV}Z_V(T,\{N_s^{(m)}\})\dd V.
\]
Density is therefore an output:
\[
\rho_{\rm mass}=
\frac{\sum_sN_s^{(m)}M_s}{N_A V_m^\star}.
\]

\section{Molecular energy operator}

One interaction form applies to every recipe:
\[
\boxed{\begin{aligned}
U_\Theta={}&
U_{\rm bond}+U_{\rm angle}+U_{\rm torsion}+U_{\rm improper}\\
&+U_{\rm rep}+U_{\rm disp}+U_{\rm elec}+U_{\rm pol}+U_{\rm mb}.
\end{aligned}}
\]
The bonded terms are
\[
U_{\rm bond}=\sum_b\frac{k_b}{2}(r_b-r_b^0)^2,\qquad
U_{\rm angle}=\sum_a\frac{k_a}{2}(\theta_a-\theta_a^0)^2,
\]
\[
U_{\rm torsion}=\sum_d\sum_k
k_{dk}\left[1+\cos(n_{dk}\varphi_d-\delta_{dk})\right],
\]
with one declared improper functional form. For the Lennard--Jones realization,
\[
U_{\rm rep}+U_{\rm disp}=
\sum_{i<j}^{\rm nonbonded}4\epsilon_{ij}
\left[\left(\frac{\sigma_{ij}}{r_{ij}}\right)^{12}
-\left(\frac{\sigma_{ij}}{r_{ij}}\right)^6\right],
\]
\[
\sigma_{ij}=\frac{\sigma_i+\sigma_j}{2},\qquad
\epsilon_{ij}=\sqrt{\epsilon_i\epsilon_j}.
\]

All periodic electrostatic blocks use the same zero-average Green function:
\[
-\nabla^2G_H(x)=4\pi\left[
\sum_{z\in\mathbb Z^3}\delta(x-Hz)-V^{-1}\right].
\]
Permanent electrostatics are
\[
U_{\rm elec}=
\frac{1}{8\pi\varepsilon_0}\sum_{i\ne j}
q_iq_jG_H(r_i-r_j)+U_{\rm self}+U_{\rm excl}.
\]
The exclusion term applies bonded exclusions and 1--4 scaling to the same periodic
Green function, including its reciprocal-space correction.

For induced dipoles \(\mu\), define the periodic field \(E_H(q)\) and interaction
matrix \(T_H(q)\) from derivatives of \(G_H\):
\[
U_{\rm pol}(q,\mu)=
\frac12\mu^\top\alpha^{-1}\mu-\mu^\top E_H(q)
-\frac12\mu^\top T_H(q)\mu.
\]
The same declared short-range damping applies to charge--dipole and dipole--dipole
blocks. The physical induced dipoles and energy are
\[
\mu^\star(q)=\arg\min_\mu U_{\rm pol}(q,\mu),\qquad
U_{\rm pol}(q)=U_{\rm pol}(q,\mu^\star(q)).
\]
The many-body term \(U_{\rm mb}\) is a declared permutation-, rotation-, and
translation-invariant function of atomic geometry. Forces and virial derive from this
single energy:
\[
F_i=-\nabla_{r_i}U_\Theta,\qquad
\Xi=-\frac{\partial U_\Theta}{\partial H}H^\top.
\]

\section{Equilibrium measure}

At \(V_m^\star\),
\[
\boxed{
\rho_\Theta(q)\dd q=
Z^{-1}e^{-\beta U_\Theta(q;V_m^\star)}\dd q.}
\]
Rigid constraints \(g(q)=0\) restrict this measure to the constraint manifold with its
induced volume element. Equilibrium integration uses a sequence \(\mathcal Q_r\) with
\[
\mathcal Q_r[f]\longrightarrow
\langle f\rangle_{\rho_\Theta}
=\int f(q)\rho_\Theta(q)\dd q.
\]
Pair basins, coordination shells, aggregates, cages, and ligand states are smooth
observables of \(q\). They do not alter \(\rho_\Theta\).

\section{Molecular coordinates and null modes}

For molecule \(a\),
\[
R_a=\frac{1}{M_a}\sum_{i\in a}m_ir_i,\quad
V_a=\frac{1}{M_a}\sum_{i\in a}m_iv_i,\quad
F_a=\sum_{i\in a}F_i.
\]
Collect the molecular translational and rotational velocities in \(V\), with
mass--inertia matrix \(M\). Let the columns of \(B\) span uniform translation and exact
rigid-constraint null modes. The physical-range projector is
\[
\boxed{\Pi=I-B(B^\top B)\pinv B^\top.}
\]
Memory, friction, diffusion, and Poisson operators act on \(\operatorname{ran}\Pi\).

\section{Microscopic dynamics for molecular memory}

The memory kernel is generated before the configuration diffusion and therefore
cannot be evolved by \(L_\Theta\).  Introduce constrained atomic momenta \(p\), mass
matrix \(M_{\rm at}\), and the phase-space Hamiltonian
\[
\mathcal H_\Theta(q,p)=
\frac12p^\top M_{\rm at}^{-1}p+U_\Theta(q).
\]
Write \(G(q)=\nabla_qg(q)\), \(v=M_{\rm at}^{-1}p\), and
\(A_g(q)=G(q)M_{\rm at}^{-1}G(q)^\top\).  On the cotangent bundle of the
molecular constraint manifold,
\[
\dot q=v,\qquad
\dot p=-\nabla_qU_\Theta-G(q)^\top\lambda,
\]
where the acceleration constraint determines the multiplier without an additional
physical parameter:
\[
\lambda=A_g(q)\pinv
\left[\dot G(q;v)v-G(q)M_{\rm at}^{-1}\nabla_qU_\Theta\right].
\]
The microscopic generator is therefore
\[
\boxed{
\mathcal G_\Theta f=
(M_{\rm at}^{-1}p)^\top\nabla_q f
-\left(\nabla_qU_\Theta+G^\top\lambda\right)^\top\nabla_p f,}
\]
on \(g(q)=0\) and \(G(q)M_{\rm at}^{-1}p=0\).  Its invariant canonical measure is
\[
\boxed{
\varrho_\Theta(q,p)\dd q\dd p=
\mathcal Z^{-1}e^{-\beta\mathcal H_\Theta(q,p)}
\dd\lambda_g(q)\dd\lambda_{T_q}(p).}
\]
Here \(\lambda_g\) and \(\lambda_{T_q}\) are the induced configuration and tangent
momentum measures.  This phase-space dynamics uses exactly the energy, masses,
cell, and constraints that define \(\rho_\Theta\).  It is used to evaluate molecular
memory; it is not the later homogenized configuration process.

\section{Conditional memory and diffusion}

Let \(\mathcal P\) be conditional expectation in
\(L^2(\varrho_\Theta)\) onto the resolved molecular centers, orientations, and their
conjugate translational and rotational momenta.  Set
\(\mathcal Q=I-\mathcal P\).  The orthogonal molecular force and its orthogonal
evolution are
\[
F^\perp=\mathcal QF,\qquad
F^\perp(t)=e^{t\mathcal Q\mathcal G_\Theta\mathcal Q}F^\perp.
\]
The fluctuation--dissipation definition of the conditional memory kernel is
\[
\boxed{
K_\Theta(q,t)=\beta\,
\mathbb E_{\varrho_\Theta}\!\left[
\Pi F^\perp(t)F^\perp(0)^\top\Pi\mid q(0)=q
\right].}
\]
The conditional expectation integrates canonical momenta and unresolved molecular
coordinates at fixed resolved configuration.  Thus \(K_\Theta\) is determined by
\((U_\Theta,M_{\rm at},\varrho_\Theta,\mathcal P)\), independently of
\(D_\Theta\) and \(L_\Theta\).
Define
\[
C_{VV}(q,t)=
\mathbb E[\Pi V(t)V(0)^\top\Pi\mid q(0)=q],
\qquad C_{PV}(q,t)=MC_{VV}(q,t).
\]
These obey the matrix Volterra equation
\[
\frac{\partial C_{PV}(q,t)}{\partial t}
=-\int_0^tK_\Theta(q,\tau)C_{VV}(q,t-\tau)\dd\tau.
\]
For \(s>0\),
\[
\boxed{
\widetilde K_\Theta(q,s)=
\left[C_{PV}(q,0)-s\widetilde C_{PV}(q,s)\right]
\widetilde C_{VV}(q,s)\pinv.}
\]
The zero-frequency friction and molecular diffusion are
\[
\boxed{\begin{aligned}
K_{\Theta,0}(q)&=\lim_{s\downarrow0}
\Pi\widetilde K_\Theta(q,s)\Pi,\\
D_\Theta(q)&=k_BT\,\Pi K_{\Theta,0}(q)\pinv\Pi.
\end{aligned}}
\]
For finite trajectory resolution \(\ell\), evaluate the same zero-frequency
operator through the positive covariance-growth form
\[
\boxed{
D_{\Theta,\ell}(q)=\frac{1}{2t_\ell}
\mathbb E\!\left[
\Delta R_\ell\Delta R_\ell^\top\mid q(0)=q
\right],\qquad
K_{\Theta,0,\ell}=k_BT\,D_{\Theta,\ell}\pinv,}
\]
where
\(
\Delta R_\ell=\int_0^{t_\ell}\Pi V(t)\dd t
\)
and the conditional mean displacement is removed.  This form is positive
semidefinite at every finite resolution and converges to the same zero-frequency
mobility as the Laplace--Volterra expression when the velocity correlation is
integrable.  It avoids replacing a noisy indefinite kernel estimate by clipped
eigenvalues.
The limit is evaluated on the physical range. Conserved and explicitly resolved slow
coordinates remain outside the inversion.

The replacement of the non-Markovian molecular equation by \(D_\Theta(q)\) is its
zero-frequency homogenization.  When the orthogonal-force kernel is integrable and
the eliminated variables mix on a finite time scale, this homogenized operator has
the same zero-frequency molecular mobility as the generalized Langevin equation.

\section{Constrained reversible generator}

Let \(T_q\) project Cartesian molecular increments onto the tangent space of the rigid
constraints. Replace \(D_\Theta\) by \(T_qD_\Theta T_q^\top\) and retain the same
symbol. The production generator is
\[
\boxed{
L_\Theta f=\rho_\Theta^{-1}\nabla_q\!\cdot
\left(\rho_\Theta D_\Theta\nabla_qf\right).}
\]
In coordinates,
\[
L_\Theta f=
-\beta\nabla U_\Theta^\top D_\Theta\nabla f
+(\nabla\!\cdot D_\Theta)^\top\nabla f
+D_\Theta:\nabla^2f.
\]
The associated stochastic dynamics is
\[
\dd q_t=
\left[-\beta D_\Theta\nabla U_\Theta
+\nabla\!\cdot D_\Theta\right]\dd t
+\sqrt{2D_\Theta}\dd W_t.
\]
The divergence drift and correlated noise belong to the same operator. For smooth
constraint-compatible functions,
\[
-\langle f,L_\Theta g\rangle_{\rho_\Theta}
=\left\langle\nabla f^\top D_\Theta\nabla g\right\rangle_{\rho_\Theta}.
\]
Thus \(L_\Theta\) is reversible and negative semidefinite on mean-zero functions.

\section{Charge polarization}

Use formal molecular charges and unwrapped molecular centers:
\[
\boxed{P(q)=\sum_aQ_aR_a^{\rm unwrap}.}
\]
Molecular translation carries unbounded charge displacement. Intramolecular rotation,
vibration, and induced polarization are bounded unless molecular charge topology
changes. The drift current is
\[
j_i(q)=L_\Theta P_i(q).
\]
The stochastic quadratic variation supplies
\[
\boxed{
C^{\rm dir}_{ij}=
\left\langle\nabla P_i^\top D_\Theta\nabla P_j
\right\rangle_{\rho_\Theta}.}
\]

\section{Poisson corrector and scalar conductivity}

For each Cartesian component solve
\[
\boxed{
-L_\Theta\chi_i=L_\Theta P_i,\qquad
\langle\chi_i\rangle_{\rho_\Theta}=0.}
\]
The weak equation is
\[
\left\langle\nabla\varphi^\top D_\Theta\nabla\chi_i\right\rangle
=-\left\langle\nabla\varphi^\top D_\Theta\nabla P_i\right\rangle
\]
for every admissible mean-zero test function \(\varphi\). The homogenized charge
diffusion is
\[
\boxed{
\mathcal D^Q_{ij}=C^{\rm dir}_{ij}
-\left\langle L_\Theta P_i,
(-L_\Theta)\pinv L_\Theta P_j\right\rangle_{\rho_\Theta}.}
\]
Equivalently,
\[
\mathcal D^Q_{ii}=
\min_{\chi\in\mathcal H_E}
\left\langle(\nabla P_i+\nabla\chi)^\top
D_\Theta(\nabla P_i+\nabla\chi)\right\rangle_{\rho_\Theta}.
\]
With \(P\) measured in \(\mathrm{C\,m}\),
\[
\boxed{
\sigma(\mathcal R,T,p)=
\frac{1}{3k_BTV_m^\star}\tr\mathcal D^Q
\quad[\mathrm{S\,m^{-1}}].}
\]

\section{Constructive nested basis}

Generate smooth permutation-, translation-, and rotation-compatible functions from
radial atomic densities, angular invariants, molecular orientations, coordination
numbers, charge-density Fourier modes, concentration modes, cluster graphs, and
hierarchical products. Let \(\mathcal G_n\) include radial and angular order, reciprocal
wavevector norm, correlation order, and cluster depth at most \(n\). After centering
and removal of exact null functions,
\[
V_n=\operatorname{span}\mathcal G_n,\qquad V_n\subset V_{n+1}.
\]
The hierarchy obeys
\[
\overline{\bigcup_{n=1}^\infty V_n}^{\|\cdot\|_E}=\mathcal H_E,\qquad
\|f\|_E^2=\langle\nabla f^\top D_\Theta\nabla f\rangle.
\]

For basis \(\Phi_n=(\phi_1,\ldots,\phi_{d_n})\), compute
\[
A_{\mu\nu}^{(n)}=
\left\langle\nabla\phi_\mu^\top D_\Theta\nabla\phi_\nu\right\rangle,
\]
\[
h_{\mu i}^{(n)}=
\left\langle\nabla\phi_\mu^\top D_\Theta\nabla P_i\right\rangle.
\]
The coefficient vector solves
\[
\boxed{A^{(n)}c_i^{(n)}=-h_i^{(n)}.}
\]
The finite charge diffusion and conductivity are
\[
\boxed{
\mathcal D^{Q,n}_{ij}=C^{\rm dir}_{ij}
-(h_i^{(n)})^\top(A^{(n)})\pinv h_j^{(n)},}
\]
\[
\boxed{
\sigma_{m,r,\ell,n}=
\frac{\tr[C^{\rm dir}-h_n^\top A_n\pinv h_n]}
{3k_BTV_m^\star}.}
\]

\section{Residual-driven enrichment}

For \(g\in V_{n+1}\setminus V_n\), define
\[
\widetilde A_{gg}
=A_{gg}-A_{g\Phi}A_{\Phi\Phi}\pinv A_{\Phi g},
\]
\[
\widetilde h_{g,i}
=h_{g,i}-A_{g\Phi}A_{\Phi\Phi}\pinv h_{\Phi,i}.
\]
When \(\widetilde A_{gg}>0\), the exact gain in the variational correction is
\[
S(g)=\sum_{i=1}^3
\frac{|\widetilde h_{g,i}|^2}{\widetilde A_{gg}}.
\]
Candidates come from the next declared hierarchy level. Exhausting one finite
dictionary does not establish convergence.

\section{Galerkin convergence}

Let \(\chi_i^\star\) solve the full Poisson problem and \(\chi_{n,i}\) be its Galerkin
projection. Galerkin orthogonality gives
\[
\langle\nabla(\chi_i^\star-\chi_{n,i})^\top
D_\Theta\nabla\phi\rangle=0,\qquad \phi\in V_n.
\]
Expanding the variational energy gives
\[
\boxed{
\sigma_{m,r,\ell,n}-\sigma_{m,r,\ell,\infty}
=\frac{1}{3k_BTV_m^\star}
\sum_{i=1}^3\|\chi_i^\star-\chi_{n,i}\|_E^2\ge0.}
\]
Therefore
\[
\sigma_{m,r,\ell,n}\downarrow\sigma_{m,r,\ell,\infty}.
\]
When equilibrium integration converges, the zero-frequency memory limit exists on the
physical range, the basis hierarchy is energy-norm dense, and the thermodynamic limit
exists,
\[
\boxed{
\sigma(\mathcal R,T,p)=
\lim_{m\to\infty}\lim_{r\to\infty}
\lim_{\ell\to\infty}\lim_{n\to\infty}
\sigma_{m,r,\ell,n}.}
\]
Every limit refines a named numerical approximation while
\(U_\Theta,K_\Theta,P,T,p\), and the intensive recipe remain fixed.

\section{Executable calculation order}

\begin{enumerate}[leftmargin=*,nosep]
\item Decode \(\mathcal R\) and construct neutral integer counts.
\item Evaluate \(U_\Theta\) and equilibrate at \((T,p)\) to obtain \(V_m^\star\)
      and samples of \(\rho_\Theta\).
\item Construct \(\mathcal G_\Theta\) from the same energy, masses, cell, and
      molecular constraints; sample its canonical phase-space measure.
\item Project atomic coordinates, momenta, velocities, and forces to constrained
      molecular variables and form \(\mathcal P\) and \(\mathcal Q\).
\item Evolve orthogonal forces with
      \(e^{t\mathcal Q\mathcal G_\Theta\mathcal Q}\), evaluate their conditional
      correlations, and solve for
      \(K_{\Theta,0}(q)\) on the physical range.
\item Form \(D_\Theta(q)=k_BT K_{\Theta,0}(q)\pinv\).
\item Evaluate \(C^{\rm dir}\), \(A_n\), and \(h_n\) over the same equilibrium measure.
\item Solve \(A_nc_{n,i}=-h_{n,i}\) after removing exact null modes.
\item Return
\[
\boxed{
\sigma_{m,r,\ell,n}=
\frac{\tr(C^{\rm dir}-h_n^\top A_n\pinv h_n)}
{3k_BTV_m^\star}.}
\]
\item Refine \((m,r,\ell,n)\) without changing the physical model.
\end{enumerate}

\section{Model boundary}

The equations contain no configured SSIP or CIP fractions, mass-action population
allocation, species conductivity bonus, terminal conductivity multiplier, or
recipe-specific transport branch. Species data determine molecular construction and
interaction parameters. The equilibrium measure determines structure. Conditional
memory determines diffusion. The Poisson corrector determines correlated charge
transport. Their composition returns the scalar conductivity.

\end{multicols}
\end{document}
